\documentclass[10pt,a4paper]{amsart}
\usepackage[margin=19mm]{geometry}
\usepackage{amsmath,amssymb,mathtools}
\usepackage[scr=rsfs,cal=euler]{mathalfa}
\usepackage{xfrac}
\usepackage[unicode,hidelinks,bookmarks=false]{hyperref}
\newcommand{\ie}{i.e.\ }
\newcommand{\cf}{cf.\ }
\newcommand{\resp}{respectively\ }
\newcommand{\Subsection}{\subsection}
\providecommand{\nobreakdash}{\nobreak\hbox{-}}
\setlength{\emergencystretch}{2em}
\multlinegap0pt
\usepackage{rotating}
\usepackage{xy}
\usepackage{amscd,amsopn,graphics,verbatim,calc}
\numberwithin{equation}{section}
\def\thefootnote{*}
\renewcommand{\subjclassname}{%
	\textup{2020} Mathematics Subject Classification}
\newtheorem{theorem}{Theorem}[section]
\newtheorem{lemma}[theorem]{Lemma}
\newtheorem{proposition}[theorem]{Proposition}
\newtheorem{corollary}[theorem]{Corollary}
\newtheorem{notation}[theorem]{Notation}
\newtheorem{problem}[theorem]{Problem}
\newtheorem{observation}[theorem]{Observation}
\newcommand{\up}[1]{{{}^{#1}\!}}
\newtheorem{claim}[theorem]{Claim}
\newtheorem{definition}[theorem]{Definition}
\newtheorem{remark}[theorem]{Remark}
\newtheorem{fact}[theorem]{Fact}
\newtheorem{example}[theorem]{Example}
\newtheorem{question}[theorem]{Question}\newtheorem{discussion}[theorem]{Discussion}
\newtheorem{conjecture}[theorem]{Conjecture}
\newtheorem{answer}[theorem]{Answer}
\newtheorem{acknowledgement}{Acknowledgement}
\renewcommand{\theacknowledgement}{}
\newcommand{\wtd}{\widetilde}
\newcommand{\Ass}{\operatorname{Ass}}
\newcommand{\Syz}{\operatorname{Syz}}
\newcommand{\grade}{\operatorname{grade}}
\newcommand{\Min}{\operatorname{Min}}
\newcommand{\red}{\operatorname{red}}
\newcommand{\Spec}{\operatorname{Spec}}
\newcommand{\spec}{\operatorname{spec}}
\newcommand{\h}{\operatorname{h}}
\newcommand{\Gpd}{\operatorname{Gpd}}
\newcommand{\Tr}{\operatorname{Tr}}
\newcommand{\rad}{\operatorname{rad}}
\newcommand{\cd}{\operatorname{cd}}
\newcommand{\Char}{\operatorname{Char}}
\newcommand{\End}{\operatorname{End}}
\newcommand{\Ht}{\operatorname{ht}}
\newcommand{\id}{\operatorname{id}}
\newcommand{\Cl}{\operatorname{Cl}}
\newcommand{\Rfd}{\operatorname{Rfd}}
\newcommand{\Gid}{\operatorname{Gid}}
\newcommand{\pd}{\operatorname{pd}}
\newcommand{\Gdim}{\operatorname{Gdim}}
\newcommand{\gdim}{\operatorname{gd}}
\newcommand{\Gfd}{\operatorname{Gfd}}
\newcommand{\E}{\operatorname{E}}
\newcommand{\Var}{\operatorname{Var}}
\newcommand{\injdim}{\operatorname{injdim}}
\newcommand{\Ext}{\operatorname{Ext}}
\newcommand{\Supp}{\operatorname{Supp}}
\newcommand{\Cone}{\operatorname{Cone}}
\newcommand{\Tor}{\operatorname{Tor}}
\newcommand{\Hom}{\operatorname{Hom}}
\newcommand{\RHom}{\operatorname{RHom}}
\newcommand{\Zd}{\operatorname{Zd}}
\newcommand{\Ann}{\operatorname{Ann}}
\newcommand{\Rad}{\operatorname{Rad}}
\newcommand{\wdim}{\operatorname{w.dim}}
\newcommand{\Ndim}{\operatorname{Ndim}}
\newcommand{\F}{\operatorname{F}}
\newcommand{\depth}{\operatorname{depth}}
\newcommand{\pdepth}{\operatorname{p.depth}}
\newcommand{\Coass}{\operatorname{Coass}}
\newcommand{\Ker}{\operatorname{Ker}}
\newcommand{\Coker}{\operatorname{Coker}}
\newcommand{\im}{\operatorname{im}}
\newcommand{\Cosupp}{\operatorname{Cosupp}}
\newcommand{\Max}{\operatorname{Max}}
\newcommand{\ux}{\underline{x}}
\newcommand{\bac}{\backslash}
\newcommand{\al}{\alpha}
\newcommand{\BN}{\mathbb{N}}
\newcommand{\BZ}{\mathbb{Z}}
\newcommand{\vpl}{\operatornamewithlimits{\varprojlim}}
\newcommand{\vli}{\operatornamewithlimits{\varliminf}}
\newcommand{\vil}{\operatornamewithlimits{\varinjlim}}
\newcommand{\bcu}{\bigcup\limits}
\newcommand{\bca}{\bigcap\limits}
\newcommand{\botm}{\bigotimes\limits}
\newcommand{\bopm}{\bigoplus\limits}
\newcommand{\uy}{\underline{y}}
\newcommand{\lo}{\longrightarrow}
\newcommand{\fm}{\mathfrak{m}}
\newcommand{\fp}{\mathfrak{p}}
\newcommand{\fq}{\mathfrak{q}}
\newcommand{\fa}{\mathfrak{a}}
\newcommand{\fb}{\mathfrak{b}}
\newcommand{\fc}{\mathfrak{c}}
\newcommand{\fx}{\underline{x}}
\newcommand{\fn}{\mathfrak{n}}
\newcommand{\fM}{\mathfrak{M}}
\newcommand{\fI}{\mathfrak{I}}
\newcommand{\fU}{\mathfrak{U}}
\newcommand{\fv}{\mathfrak{v}}
\newcommand{\fX}{\mathfrak{X}}
\newcommand{\fW}{\mathfrak{W}}
\newcommand{\fQ}{\mathfrak{Q}}
\newcommand{\fP}{\mathfrak{P}}
\newcommand{\fZ}{\mathfrak{Z}}
\newcommand{\one}{\ensuremath{\mathord \ast}}
\newcommand{\two}{\ensuremath{\mathord \dagger}}
\newcommand{\CMdef}{\operatorname{CMdef}}
\newcommand{\fC}{\mathfrak{C}}
\newcommand{\uc}{\underline{c}}
\DeclareMathOperator{\Frac}{Frac}
\def\mapdown#1{\Big\downarrow\rlap{$\vcenter{\hbox{$\scriptstyle#1$}}$}}
\renewcommand{\baselinestretch}{1.27}
\newenvironment{prf}[1][Proof]{\begin{proof}[\bf #1]}{\end{proof}}
\begin{document}
\title[Non-Noetherian Bass and Betti numbers]{Non-Noetherian Bass and Betti numbers}
\author{Mohsen Asgharzadeh; Elham Mahdavi}
\date{}
\maketitle
\noindent\emph{Source and licence.} Non-Noetherian Bass and Betti numbers. Version arXiv:2606.20391v4 (CC BY 4.0). This material is distributed under the Creative Commons Attribution 4.0 International licence, http://creativecommons.org/licenses/by/4.0/, as recorded in the arXiv OAI licence metadata for this exact version and shown on the abstract page https://arxiv.org/abs/2606.20391v4. Full rights notice: licenses/arxiv-2606.20391-CC-BY-4.0.txt.\par\smallskip
\noindent\textit{Editorial scope.} Counted unit: the original \texttt{theorem} environment \texttt{-} at source lines 460--462 together with its complete proof at lines 464--507; the delivered document keeps source lines 430--508 so that every referenced label and every citation resolves inside it. Native editable source is the paper's own concatenated TeX; the AMS wrapper is editorial formatting only. No publisher class, upstream script or unrelated artwork is redistributed.
	\section{Weakly tor-rigid}
This section is divided into two subsections:
	\subsection{Weak tor-rigidity and big Cohen--Macaulay algebras}
	
	This subsection focuses on tor-rigidity and big Cohen--Macaulay algebras, proving that Schoutens' question has a positive answer for tor-rigid algebras.
	
\begin{question}(Schoutens, \cite[End of \S 2]{sc1}).
	Suppose \(S\) is a big Cohen--Macaulay algebra over a local ring \(R\). If \(\operatorname{Tor}_n^R(k,S)=0\) for some \(n\ge 1\), does it follow that \(R\) is Cohen--Macaulay?
\end{question}First, recall that vanishing of all Betti numbers does not imply flatness:
\begin{example}[Bartijn--Strooker \cite{BS}]
	Let \(R = k[[x,y,z]]\) and let \(\varphi\) be a free module of infinite countable rank. Define \(\Omega := \varphi + (x,y)\widehat{\varphi}\). Then \(\{x,y,z\}\) is an \(\Omega\)-sequence, \(\{z,x,y\}\) is not an \(\Omega\)-sequence, and \(\operatorname{Tor}_i^R(k,\Omega)=0\).  
\end{example}
The following extends Example 3.2, and simplifies a result of Huneke (see \cite[9.1]{Hun}).

\begin{observation}\label{3.3}
	If \(R\) is a regular local ring and \(M\) is a big Cohen--Macaulay module, then \(\beta_i(M)=0\) for all \(i>0\).
\end{observation}

\begin{proof}
	Since \(R\) is regular, the residue field \(k\) admits a finite free resolution
	\[
	0 \longrightarrow F_d \longrightarrow F_{d-1} \longrightarrow \cdots \longrightarrow F_0 \longrightarrow k \longrightarrow 0.
	\]
	Tensoring this resolution with \(M\) yields a complex whose homology is precisely \(\operatorname{Tor}_i^R(k,M)\). Because \(M\) is big Cohen--Macaulay, the Koszul complex on any system of parameters is exact on \(M\), which is equivalent to the vanishing of all positive \(\operatorname{Tor}\)'s with the residue field; hence \(\operatorname{Tor}_i^R(k,M)=0\) for every \(i>0\). Consequently, the Betti numbers \(\beta_i(M)=\dim_k \operatorname{Tor}_i^R(k,M)\) vanish for all \(i>0\).
\end{proof}

\begin{definition}
	An \(R\)-module \(S\) is called \emph{weakly tor-rigid} if \(\operatorname{Tor}_j^R(R/I,S)=0\) for some parameter ideal \(I\) and some integer \(j\) implies \(\operatorname{Tor}_{j+1}^R(R/I,S)=0\).
\end{definition}
Concerning Question 3.1, he proved the desired claim when $n\leq 2$. Here, we deal with the case \(n=3\), and the general case seems being analogous.

\begin{theorem}
Suppose \(S\) is a big Cohen--Macaulay, weakly tor-rigid and  \(\operatorname{Tor}_3^R(S,k)=0\).   Then \(R\) is Cohen--Macaulay.
\end{theorem}

\begin{proof}
 There exists an exact sequence
$
	0 \to M \to F_1 \xrightarrow{\varphi} F_0 \to  S \to 0,
$
	where \(M=\Syz_2(S)\) and \(F_0,F_1\) are free modules (not necessarily of finite rank). As observed in \cite[2.5]{sc1}, it suffices to prove that \(IS\cap R=I\) for some parameter ideal \(I\). 
Consider the two short exact sequences arising from the free resolution:
	\[
	0 \longrightarrow M \longrightarrow F_1 \longrightarrow E:=\operatorname{Im}(\varphi) \longrightarrow 0,
	\]
	\[
	0 \longrightarrow E \longrightarrow F_0 \longrightarrow S \longrightarrow 0.
	\]
	Tensoring these with \(\overline{R}:=R/I\) yields the exact sequences

\begin{itemize}
	\item[(i)]
 \(0 \longrightarrow \operatorname{Tor}_1^R(E,\overline{R}) \longrightarrow M/IM \longrightarrow F_1\otimes_R \overline{R} \longrightarrow E\otimes_R \overline{R} \longrightarrow 0,\) and
 \item[(ii)]	
	\(	0 \longrightarrow \operatorname{Tor}_1^R(S,\overline{R}) \longrightarrow E\otimes_R \overline{R} \longrightarrow F_0\otimes_R \overline{R} \longrightarrow S\otimes_R \overline{R} \longrightarrow 0.
	\) 
\end{itemize}

Now, tensoring the short exact sequence \(0\to I\to R\to R/I\to 0\) with \(S\) to obtain
$
	\operatorname{Tor}_1^R(R/I,S) \to S\otimes_R I \to S.
$
	Since \(S\) is big Cohen-Macaulay, $I$ is generated by an $S$-sequence. Hence, the map \(S\otimes_R I\to S\) is injective, so \(\operatorname{Tor}_1^R(R/I,S)=0\). Moreover, the natural map \(I\otimes_R S\to IS\) is an isomorphism. Applying tor-rigidity, we also get \(\operatorname{Tor}_2^R(R/I,S)=0\). Assuming \(\operatorname{Tor}_1^R(R/I,E)=\operatorname{Tor}_2^R(R/I,S)=0\), the first tensored sequence reduces to
	\[
\zeta:=	0\longrightarrow M/IM \longrightarrow F_1\otimes_R \overline{R} \longrightarrow E\otimes_R \overline{R} \longrightarrow 0.
	\]
	From this we deduce \(\operatorname{Tor}_1^R(M,k)=\operatorname{Tor}_3^R(S,k)=0\), and consequently \(\operatorname{Tor}_1^{\overline{R}}(M/IM,k)=0\) by \cite[2.1]{sc1}. Hence \(M/IM\) is flat over \(\overline{R}\), (see \cite[22.3]{mat}). Now, by $\zeta$	we deduce \(E\otimes_R \overline{R}\) has finite flat dimension over \(\overline{R}\).
	Since \(\overline{R}\) is zero-dimensional, the flat dimension of \(E\otimes_R \overline{R}\) over \(\overline{R}\) is at most dimension of the ring, then it is actually flat over the artinian ring \(\overline{R}\). 
By 		\[
0 = \operatorname{Tor}_1^R(S,\overline{R}) \longrightarrow E\otimes_R \overline{R} \longrightarrow F_0\otimes_R \overline{R} \longrightarrow S\otimes_R \overline{R} \longrightarrow 0.
\]
we see $S\otimes_R \overline{R}$
is actually flat over   \(\overline{R}\).  Recall that flat modules are torsion-free, thus, annihilator
of $S\otimes_R \overline{R}=S/IS$
as an $R$-module is $I$.
In other words, \(IS\cap R=I\). This ensures that \(I\) is generated by an \(R\)-sequence, as
\(I\) is $S$-regular sequence.
By definition, \(R\) is Cohen--Macaulay.
\end{proof}


\begin{thebibliography}{99}
\bibitem[BS]{BS} J. Bartijn and J. Strooker, \emph{Modifications monomiales}. In M.-P. Malliavin, editor, \emph{S\'eminaire d'Alg\`ebre Paul Dubreil et Marie-Paule Malliavin: Proceedings, Paris 1982 (35\`eme Ann\'ee)}, volume 1029 of \emph{Lecture Notes in Mathematics}, pages 192--217. Springer-Verlag, Berlin, Heidelberg, 1983.
\bibitem[Hun]{Hun} {C.~Huneke}, \emph{Tight closure and its applications}, CBMS Regional Conference Series in Mathematics Volume {\bf88}, 1996.
\bibitem[mat]{mat} H. Matsumura, \emph{Commutative ring theory}, Cambridge Studies in Advanced Math, \textbf{8}, (1986).
\bibitem[sc1]{sc1} H. Schoutens, \emph{On the vanishing of Tor of the absolute integral closure}, J. Algebra {\bf 275} (2004), 567–574.
\end{thebibliography}
\end{document}
